\documentclass[preview]{standalone}

\usepackage{amsmath}
\usepackage{amssymb}
\usepackage{stellar}
\usepackage{definitions}

\begin{document}

\id{ring-substructures}
\genpage

\section{Subrings and ideals}

\begin{snippetdefinition}{subring-definition}{Subring}
    Let \(A\) be a \ring. A \emph{subring} of \(A\) is a subset
    \(B\subseteq A\) that is a \ring under the operations induced by \(A\)
    and contains the same multiplicative identity. Equivalently,
    \(B\leq A\) if
    \[
        1_A\in B,
        \qquad
        b_1-b_2\in B,
        \qquad
        b_1b_2\in B
    \]
    for all \(b_1,b_2\in B\).
\end{snippetdefinition}

\begin{snippetdefinition}{ring-ideal-definition}{Ideal}
    Let \(A\) be a \ring. A subset \(I\subseteq A\) is a
    \emph{left ideal} if \((I,+)\) is a \subgroup of \((A,+)\) and
    \[
        ax\in I
    \]
    for every \(a\in A\) and \(x\in I\). It is a \emph{right ideal} if
    \[
        xa\in I
    \]
    for every \(a\in A\) and \(x\in I\). A \emph{two-sided ideal}, or
    simply an \emph{ideal}, is both a left and a right ideal; this is denoted
    by \(I\unlhd A\).
\end{snippetdefinition}

\begin{plainhtml}
    In a commutative ring, the notions of left, right, and two-sided ideal coincide.
\end{plainhtml}

\begin{snippetproposition}{ring-ideal-identity}{An ideal containing the identity}
    Let \(A\) be a \ring and let \(I\) be a left or right \ideal of \(A\).
    If \(1_A\in I\), then \(I=A\).
\end{snippetproposition}

\begin{snippetproof}{ring-ideal-identity-proof}{ring-ideal-identity}{An ideal containing the identity}
    If \(I\) is a left ideal, then \(a=a1_A\in I\) for every \(a\in A\).
    If \(I\) is a right ideal, then \(a=1_Aa\in I\) for every \(a\in A\).
\end{snippetproof}

\begin{snippetproposition}{trivial-ring-ideals}{Trivial ideals}
    Every \ring \(A\) has the two-sided \ideal[ideals]
    \[
        \{0_A\}\idealin A,
        \qquad
        A\idealin A.
    \]
\end{snippetproposition}

\begin{snippet}{zero-ideal-not-unital-subring}
    Under the convention that subrings contain the same identity as the
    ambient \ring, the zero \ideal \(\{0_A\}\) is not a \subring unless
    \(A\) is the trivial ring.
\end{snippet}

\section{Principal ideals}

\includesnpt{principal-ideal-definition}

\begin{snippetproposition}{integers-subrings-and-ideals}{Subrings and ideals of the integers}
    Under the unital-subring convention, the only \subring of \(\integers\)
    is \(\integers\) itself. Its \ideal[ideals] are exactly
    \[
        n\integers=\{kn\suchthat k\in\integers\},
        \qquad n\in\naturalnumbers\union\{0\}.
    \]
\end{snippetproposition}

\begin{snippetproof}{integers-subrings-and-ideals-proof}{integers-subrings-and-ideals}{Subrings and ideals of the integers}
    A \subring \(B\) of \(\integers\) contains \(1\), hence every positive
    integer by closure under addition and every negative integer by closure
    under additive inverses. Thus \(B=\integers\).

    For every \(n\), the \set \(n\integers\) is an \ideal. Conversely, let
    \(I\neq\{0\}\) be an \ideal of \(\integers\), and let \(n\) be its
    least positive element. Then \(n\integers\subseteq I\). For \(x\in I\),
    Euclidean division gives
    \[
        x=qn+r,
        \qquad
        0\leq r<n.
    \]
    Since \(r=x-qn\in I\), minimality forces \(r=0\). Hence
    \(I=n\integers\). The zero ideal is \(0\integers\).
\end{snippetproof}

\begin{snippettheorem}{field-polynomial-ring-ideals-principal}{Ideals of a polynomial ring over a field}
    Let \(K\) be a \field. Every \ideal of \(K[x]\) is principal.
\end{snippettheorem}

\begin{snippetproof}{field-polynomial-ring-ideals-principal-proof}{field-polynomial-ring-ideals-principal}{Ideals of a polynomial ring over a field}
    Let \(I\neq\{0\}\) be an \ideal of \(K[x]\), and choose
    \(f\in I\setminus\{0\}\) of minimal degree. Since \(f\in I\),
    \((f)\subseteq I\). For \(g\in I\), polynomial division gives
    \[
        g=qf+r,
        \qquad
        r=0\text{ or }\deg r<\deg f.
    \]
    Since \(r=g-qf\in I\), minimality of \(\deg f\) forces \(r=0\).
    Therefore \(g\in(f)\), so \(I=(f)\). The zero ideal is generated by
    the zero polynomial.
\end{snippetproof}

\begin{snippetproposition}{field-polynomial-principal-generator-uniqueness}{Minimal-degree generators differ by a scalar}
    Let \(K\) be a \field, let \(I=(f)\idealin K[x]\) with \(f\neq0\), and
    let \(g\in I\) satisfy \(\deg g=\deg f\). Then
    \[
        g=\lambda f
    \]
    for some \(\lambda\in K^\times\).
\end{snippetproposition}

\begin{snippetproof}{field-polynomial-principal-generator-uniqueness-proof}{field-polynomial-principal-generator-uniqueness}{Minimal-degree generators differ by a scalar}
    Since \(g\in(f)\), there is \(q\in K[x]\) such that \(g=qf\). Thus
    \[
        \deg g=\deg q+\deg f.
    \]
    The equality \(\deg g=\deg f\) implies \(\deg q=0\), so
    \(q=\lambda\in K^\times\).
\end{snippetproof}

\begin{snippet}{unit-in-ideal-characterization}
    Let \(A\) be a \ring and \(I\idealin A\). If \(I\) contains a unit
    \(u\), then \(1_A=u^{-1}u\in I\), and hence \(I=A\). In particular, an
    \ideal of \(K[x]\) containing a nonzero constant polynomial is the whole
    ring.
\end{snippet}

\section{Centralizers and center}

\begin{snippetdefinition}{ring-element-centralizer-definition}{Centralizer of an element in a ring}
    Let \(A\) be a \ring and \(b\in A\). The \emph{centralizer of \(b\) in
    \(A\)} is the \set
    \[
        C_A(b)=\{a\in A\suchthat ab=ba\}.
    \]
\end{snippetdefinition}

\begin{snippetdefinition}{ring-subset-centralizer-definition}{Centralizer of a subset in a ring}
    Let \(A\) be a \ring and \(S\subseteq A\). The \emph{centralizer of
    \(S\) in \(A\)} is the \set
    \[
        C_A(S)
        =\{a\in A\suchthat as=sa\text{ for every }s\in S\}
        =\bigcap_{s\in S}C_A(s).
    \]
\end{snippetdefinition}

\begin{snippetdefinition}{ring-center-definition}{Center of a ring}
    Let \(A\) be a \ring. The \emph{center of \(A\)} is the \set
    \[
        Z(A)=C_A(A)
        =\{a\in A\suchthat ab=ba\text{ for every }b\in A\}.
    \]
\end{snippetdefinition}

\begin{snippetproposition}{ring-centralizers-center-subrings}{Centralizers and the center are subrings}
    Let \(A\) be a \ring. For every \(b\in A\) and \(S\subseteq A\),
    \[
        C_A(b)\subringleq A,
        \qquad
        C_A(S)\subringleq A,
        \qquad
        Z(A)\subringleq A.
    \]
\end{snippetproposition}

\begin{snippetproof}{ring-centralizers-center-subrings-proof}{ring-centralizers-center-subrings}{Centralizers and the center are subrings}
    Fix \(b\in A\). Both \(0_A\) and \(1_A\) commute with \(b\). If
    \(a_1,a_2\in C_A(b)\), then
    \[
        (a_1-a_2)b=b(a_1-a_2)
    \]
    and
    \[
        (a_1a_2)b=a_1(a_2b)=a_1(ba_2)=(ba_1)a_2=b(a_1a_2).
    \]
    Hence \(C_A(b)\subringleq A\). Since
    \(C_A(S)=\bigcap_{s\in S}C_A(s)\), the second assertion follows from
    closure of subrings under intersections. The third is the case \(S=A\).
\end{snippetproof}

\section{Intersections, unions, and sums}

\begin{snippetproposition}{ring-subsets-intersection}{Intersection of subrings}
    Let \(A\) be a \ring and let \(\{B_i\}_{i\in J}\) be a nonempty family
    of \subring[subrings] of \(A\). Then
    \[
        \bigcap_{i\in J}B_i\subringleq A.
    \]
\end{snippetproposition}

\begin{snippetproof}{ring-subsets-intersection-proof}{ring-subsets-intersection}{Intersection of subrings}
    Every \(B_i\) contains \(1_A\). If \(a,b\) belong to every \(B_i\),
    then \(a-b\) and \(ab\) belong to every \(B_i\). Hence the intersection
    satisfies the subring criterion.
\end{snippetproof}

\begin{snippetproposition}{subring-union-sum-need-not-subring}{Unions and sums of subrings}
    Let \(A\) be a \ring and let \(B_1,B_2\subringleq A\). In general,
    neither
    \[
        B_1\union B_2
        \qquad\text{nor}\qquad
        B_1+B_2=\{b_1+b_2\suchthat b_i\in B_i\}
    \]
    is a \subring of \(A\).
\end{snippetproposition}

\begin{snippetexample}{subring-union-sum-counterexample}{A union and a sum that are not subrings}
    In \(\complexnumbers\), let
    \[
        B_1=\integers[i],
        \qquad
        B_2=\integers[\sqrt2].
    \]
    Both are \subring[subrings]. However,
    \(i\in B_1\), \(\sqrt2\in B_2\), and
    \(i+\sqrt2\notin B_1\union B_2\), so their union is not closed under
    addition. Moreover,
    \[
        B_1+B_2
        =\{a+b\sqrt2+ci\suchthat a,b,c\in\integers\},
    \]
    which contains \(i\) and \(\sqrt2\) but not \(i\sqrt2\). Thus the sum
    is not closed under multiplication.
\end{snippetexample}

\end{document}
